% English abstract
\clearpage
\phantomsection
\addcontentsline{toc}{chapter}{Abstract}

\begin{singlespace}
\begin{latin}
\begin{center}
{\Large\bfseries Abstract}\\[0.4cm]
\end{center}

\setstretch{1.10}
The numerical solution of nonlinear partial differential equations (PDEs), particularly the viscous Burgers equation characterized by steep shock gradients, represents a fundamental challenge in scientific computing and fluid dynamics. While classical numerical methods are mathematically rigorous, they inherently rely on computational mesh generation and suffer from exponential complexity in higher dimensions. On the other hand, purely data-driven deep learning models require vast amounts of labeled training data and fail to ensure physical consistency in sparse-data regimes. The physics-informed neural network (PINN) paradigm addresses these limitations by embedding governing differential laws directly into the loss function of deep neural networks. This thesis presents a comprehensive theoretical study, computational implementation, and empirical evaluation of continuous-time PINNs for solving the forward problem of the one-dimensional nonlinear Burgers equation, benchmarked against authoritative literature.

First, the theoretical foundations spanning partial differential equations, classical numerical methods, deep feedforward networks, optimization algorithms, and automatic differentiation are established. Next, the continuous-time PINN formulation is constructed for the nonlinear Burgers equation using automatic differentiation to evaluate physical differential residuals. In the computational section, the continuous-time model is implemented in Python using the JAX library and trained via a two-stage optimization strategy combining the Adam optimizer and the quasi-Newton L-BFGS method. An independent high-precision reference solver based on the Fourier spectral method with fourth-order exponential time-differencing (ETDRK4) is developed and verified against the semi-analytical Cole-Hopf transformation.

The empirical results and extensive comparative analysis with the seminal benchmark of Raissi et al. (2019) demonstrate that the physics-informed neural network accurately reconstructs the continuous spatio-temporal solution field using only a small set of boundary and initial points alongside collocation points. The model effectively captures the sharp shock front dynamics while eliminating the need for computational mesh generation, highlighting the high data efficiency and robust approximation capability of physics-informed architectures for nonlinear dynamical systems.

\vspace{0.3cm}
\noindent\textbf{Keywords:} \KeywordsEn.
\end{latin}
\end{singlespace}
